\documentclass[12pt,a4paper]{article}

% ==============================================================================
% PAQUETES BÁSICOS Y CODIFICACIÓN
% ==============================================================================
\usepackage[T1]{fontenc}
\usepackage[utf8]{inputenc}
\usepackage[spanish,es-tabla,es-noquoting]{babel}

% Márgenes idénticos al documento Word original (2.3 cm laterales, 2.0 cm sup/inf)
\usepackage[a4paper, left=2.3cm, right=2.3cm, top=2.0cm, bottom=2.0cm]{geometry}

% Fuente Calibri idéntica a Word mediante la tipografía Carlito
\usepackage{carlito}
\renewcommand{\familydefault}{\sfdefault}

% Paquetes para tablas, gráficos y diseño
\usepackage{graphicx}
\usepackage{array}
\usepackage{tabularx}
\usepackage{booktabs}
\usepackage{amsmath,amssymb}
\usepackage{float}
\usepackage{caption}
\usepackage{subcaption}
\usepackage{xcolor}
\usepackage{listings}
\usepackage[hidelinks]{hyperref}

% Rutas de búsqueda de imágenes (soporta carpeta img/ o directorio actual)
\graphicspath{{img/}{logos/}{./}}

% ==============================================================================
% CONFIGURACIÓN DE CÓDIGO FUENTE (Para prácticas de Big Data / Programación)
% ==============================================================================
\definecolor{codegreen}{rgb}{0,0.6,0}
\definecolor{codegray}{rgb}{0.45,0.45,0.45}
\definecolor{codepurple}{rgb}{0.55,0,0.8}
\definecolor{backcolour}{rgb}{0.97,0.97,0.97}

\lstdefinestyle{estilocodigo}{
    backgroundcolor=\color{backcolour},   
    commentstyle=\color{codegreen},
    keywordstyle=\color{blue}\bfseries,
    numberstyle=\tiny\color{codegray},
    stringstyle=\color{codepurple},
    basicstyle=\ttfamily\footnotesize,
    breakatwhitespace=false,         
    breaklines=true,                 
    captionpos=b,                    
    keepspaces=true,                 
    numbers=left,                    
    numbersep=6pt,                  
    showspaces=false,                
    showstringspaces=false,
    showtabs=false,                  
    tabsize=2,
    frame=single,
    rulecolor=\color{codegray!40}
}
\lstset{style=estilocodigo}

% ==============================================================================
% DEFINICIÓN DE MACROS Y FORMATO DE ENCABEZADOS << SECCIÓN >>
% ==============================================================================
% Símbolos de llaves dobles exactas sin activar ligaduras ni shorthands
\newcommand{\llav}{\symbol{60}\symbol{60}}
\newcommand{\rlav}{\symbol{62}\symbol{62}}
\newcommand{\marcador}[1]{\llav\ #1\ \rlav}
\newcommand{\marcadornospace}[1]{\llav#1\rlav}

% Comando de sección al estilo exacto del Word original
\newcommand{\seccion}[1]{%
  \vspace{0.75cm}%
  {\centering\bfseries\fontsize{12pt}{14pt}\selectfont\marcador{\MakeUppercase{#1}}\par}%
  \vspace{0.35cm}%
}

% Comando de sección de referencias sin espacio interior (como en el Word original)
\newcommand{\seccionref}{%
  \vspace{0.75cm}%
  {\centering\bfseries\fontsize{12pt}{14pt}\selectfont\marcadornospace{REFERENCIAS}\par}%
  \vspace{0.35cm}%
}

% Compatibilidad con \section estándar
\makeatletter
\renewcommand{\section}{\@ifstar\mi@section@star\mi@section@nostar}
\newcommand{\mi@section@nostar}[1]{\seccion{#1}}
\newcommand{\mi@section@star}[1]{\seccion{#1}}
\makeatother

% ==============================================================================
% CONFIGURACIÓN DE BIBLIOGRAFÍA (Soporta .bib y thebibliography estándar)
% ==============================================================================
\makeatletter
\renewenvironment{thebibliography}[1]{%
  \seccionref
  \list{\@biblabel{\@arabic\c@enumiv}}%
       {\settowidth\labelwidth{\@biblabel{#1}}%
        \leftmargin\labelwidth
        \advance\leftmargin\labelsep
        \@openbib@code
        \usecounter{enumiv}%
        \let\p@enumiv\@empty
        \renewcommand\theenumiv{\@arabic\c@enumiv}}%
  \sloppy
  \clubpenalty4000
  \@clubpenalty \clubpenalty
  \widowpenalty4000%
  \sfcode`\.\@m}
  {\def\@noitemerr
    {\@latex@warning{Empty `thebibliography' environment}}%
  \endlist}
\makeatother

% ==============================================================================
% DATOS DEL TRABAJO / ASIGNATURA / ESTUDIANTE
% (Modificar estos campos según cada práctica)
% ==============================================================================
\newcommand{\NombreAsignatura}{<< NOMBRE DE ASIGNATURA >>}
\newcommand{\NombreTrabajo}{<< NOMBRE TRABAJO >>}
\newcommand{\CodigoEstudiante}{2022246951}
\newcommand{\ApellidosNombresEstudiante}{Medina Delgado Italo Gabriel}
\newcommand{\FechaEntrega}{} % Dejar vacío si se llena a mano o escribir fecha

% ==============================================================================
% INICIO DEL DOCUMENTO
% ==============================================================================
\begin{document}
\pagestyle{plain}

% ------------------------------------------------------------------------------
% 1. ENCABEZADO INSTITUCIONAL CON LOGOS (UCSM y EPIS)
% ------------------------------------------------------------------------------
\noindent\makebox[\textwidth][c]{%
  \begin{tabular}{@{} m{3.0cm} @{\hspace{0.25cm}} m{11.8cm} @{\hspace{0.25cm}} m{2.4cm} @{}}
    \centering\includegraphics[height=3cm, keepaspectratio]{logo_ucsm.png} &
    \centering
    {\fontsize{15pt}{18pt}\selectfont\bfseries UNIVERSIDAD CATÓLICA DE SANTA MARÍA}\par\vspace{2.5pt}
    \rule{11.6cm}{0.5pt}\par\vspace{3.5pt}
    {\fontsize{13pt}{16pt}\selectfont\bfseries ESCUELA PROFESIONAL DE INGENIERÍA DE SISTEMAS} &
    \centering\includegraphics[height=3cm, keepaspectratio]{logo_epis.png}
  \end{tabular}%
}

\vspace{0.9cm}

% ------------------------------------------------------------------------------
% 2. TÍTULOS DE ASIGNATURA Y TRABAJO
% ------------------------------------------------------------------------------
\begin{center}
  {\fontsize{12pt}{14pt}\selectfont\bfseries\NombreAsignatura}\par\vspace{4pt}
  {\fontsize{12pt}{14pt}\selectfont\bfseries\NombreTrabajo}
\end{center}

\vspace{0.4cm}

% ------------------------------------------------------------------------------
% 3. TABLA DE DATOS DEL ALUMNO
% ------------------------------------------------------------------------------
\noindent
\renewcommand{\arraystretch}{1.35}
\begin{tabularx}{\textwidth}{|>{\raggedright\arraybackslash}m{2.7cm}|>{\raggedright\arraybackslash}X|>{\centering\arraybackslash}m{2.6cm}|}
  \hline
  \centering\bfseries CÓDIGO/DNI & \centering\bfseries APELLIDOS Y NOMBRES & \bfseries FECHA \tabularnewline
  \hline
  \CodigoEstudiante & \ApellidosNombresEstudiante & \FechaEntrega \tabularnewline
  \hline
\end{tabularx}

\vspace{0.6cm}

% ------------------------------------------------------------------------------
% 4. SECCIONES DEL TRABAJO / PRÁCTICA
% ------------------------------------------------------------------------------

\seccion{ACTIVIDADES}
% Escribir o desarrollar aquí las actividades de la práctica
\vspace{1.5cm}

\seccion{EJERCICIOS}
% Escribir o desarrollar aquí los ejercicios de la práctica
\vspace{1.5cm}

\seccion{CUESTIONARIO}
% Escribir o desarrollar aquí las preguntas y respuestas del cuestionario
\vspace{1.5cm}

\seccion{CONCLUSIONES}
% Escribir o redactar aquí las conclusiones de la práctica
\vspace{1.5cm}

% ------------------------------------------------------------------------------
% 5. REFERENCIAS
% ------------------------------------------------------------------------------
% Opción A (Recomendada con .bib): Descomentar las siguientes líneas si se tiene archivo .bib:
% \bibliographystyle{IEEEtran}
% \bibliography{referencias} % busca referencias.bib
%
% Opción B (Manual): Descomentar para incluir referencias directas en el archivo:
% \begin{thebibliography}{99}
%   \bibitem{ref1} Autor, "Título del recurso o artículo", Revista/Editorial, 2026.
% \end{thebibliography}
%
% Opción C (Plantilla base vacía como en el Word original):
\seccionref

\end{document}
